\documentclass[multi=circuitikz,border=10pt]{standalone}
\usepackage[utf8]{inputenc}
\usepackage[T1]{fontenc}
\usepackage{amsmath}
\usepackage{circuitikz}
\ctikzset{inductors/scale=0.8}
\tikzset{bus/.style={line width=2.2pt}}

% ---------- Macros reutilizadas nas figuras 1 a 4 ----------
\newcommand{\refline}{\draw[thick] (-1,0) -- (11.8,0) node[right]{Ref.};}
\newcommand{\busbars}{%
  \foreach \x/\t in {3.6/{G (15 kV)}, 6.6/{T (245 kV)}, 9.6/{F}} {
    \draw[bus] (\x,2.1) -- (\x,2.9); \node[font=\scriptsize] at (\x,3.75) {\t}; }}
\newcommand{\gerador}[1]{\draw (0.6,2.5) to[L, l={#1}] (3.0,2.5) -- (3.6,2.5);}
\newcommand{\trafo}{\draw (3.6,2.5) to[L, l={$X_t = j0{,}03$ pu}] (6.0,2.5) -- (6.6,2.5);}
\newcommand{\linha}[2]{\draw (6.6,2.5) to[L, l={#1}, i_={#2}] (9.0,2.5) -- (9.6,2.5);}

\begin{document}
% =====================================================
% FIGURA 1 - Sequencia POSITIVA
% =====================================================
\begin{circuitikz}[american, font=\footnotesize]
  \refline \busbars
  \draw (0,0) to[american voltage source, v_=$E_b$, invert] (0,2.5) -- (0.6,2.5);
  \gerador{$X_{g1} = j0{,}045$ pu}
  \trafo
  \linha{$Z_{1,L} = j0{,}0666$ pu}{$I_b^{(1)}$}
  \draw[dashed, gray] (9.6,2.1) -- (9.6,0);
  \draw (9.6,2.1) to[open, v^=$V_b^{(1)}$] (9.6,0);
  \node[font=\small] at (5,-1.1) {Rede de sequência positiva: $Z_1^{tot} = j0{,}045 + j0{,}03 + j0{,}0666 = j0{,}1416$ pu};
\end{circuitikz}

% =====================================================
% FIGURA 2 - Sequencia NEGATIVA
% =====================================================
\begin{circuitikz}[american, font=\footnotesize]
  \refline \busbars
  \draw (0,0) -- (0,2.5) -- (0.6,2.5);
  \gerador{$X_{g2} = j0{,}045$ pu}
  \trafo
  \linha{$Z_{2,L} = j0{,}0666$ pu}{$I_b^{(2)}$}
  \draw[dashed, gray] (9.6,2.1) -- (9.6,0);
  \draw (9.6,2.1) to[open, v^=$V_b^{(2)}$] (9.6,0);
  \node[font=\small] at (5,-1.1) {Rede de sequência negativa: $Z_2^{tot} = j0{,}1416$ pu (sem fonte)};
\end{circuitikz}

% =====================================================
% FIGURA 3 - Sequencia ZERO (Delta-Delta = circuito aberto)
% =====================================================
\begin{circuitikz}[american, font=\footnotesize]
  \refline \busbars
  \draw (0,0) -- (0,2.5) -- (0.6,2.5);
  \gerador{$X_{g0} = j0{,}0125$ pu}
  % transformador Delta-Delta: circuito aberto para sequencia zero
  \draw (3.6,2.5) -- (4.3,2.5); \fill (4.3,2.5) circle (1.6pt);
  \draw (4.3,2.5) -- (5.3,2.95); \fill (5.6,2.5) circle (1.6pt);
  \draw (5.6,2.5) -- (6.6,2.5);
  \node[font=\scriptsize, align=center] at (5,4.4) {$\Delta$-$\Delta$: circuito aberto\\(Cap.\,11, Fig.\,11.17, caso 3)};
  \node[font=\scriptsize] at (5,1.8) {($X_t = j0{,}03$ pu não conduz $I_0$)};
  \linha{$Z_{0,L} = j0{,}2666$ pu}{$I_0 = 0$}
  \draw[dashed, gray] (9.6,2.1) -- (9.6,0);
  \draw (9.6,2.1) to[open, v^=$V_b^{(0)}$] (9.6,0);
  \node[font=\small] at (5,-1.1) {Rede de sequência zero: $Z_0$ visto de F $= \infty$ ($I_0 = 0$)};
\end{circuitikz}

% =====================================================
% FIGURA 4 - Falta TRIFASICA (so rede positiva)
% =====================================================
\begin{circuitikz}[american, font=\footnotesize]
  \refline \busbars
  \draw (0,0) to[american voltage source, v_=$E_b$, invert] (0,2.5) -- (0.6,2.5);
  \gerador{$X_{g1} = j0{,}045$ pu}
  \trafo
  \linha{$Z_{1,L} = j0{,}0666$ pu}{$I_b^{(1)}$}
  \draw[thick] (9.6,2.5) -- (9.6,0);
  \node[font=\scriptsize, align=left, anchor=west] at (9.9,1.25) {curto\\trifásico};
  \node[font=\small, align=center] at (5,-1.3) {Falta trifásica em F: só a rede positiva\\
    $I_b = E_b / Z_1^{tot} = 1{,}0 / j0{,}1416 = 7{,}06$ pu $\approx 1664$ A};
\end{circuitikz}

% =====================================================
% FIGURA 5 - Falta FASE-TERRA: redes em serie
% =====================================================
\begin{circuitikz}[american, font=\footnotesize]
  \draw (-3,7) to[american voltage source, v_={$E_b = 1\angle{-120^\circ}$ pu}, invert] (-3,9);
  % seq. positiva
  \draw (-3,9) -- (0,9) to[L, l={$Z_1^{tot} = j0{,}1416$ pu}, i_=$I_b^{(1)}$] (4.3,9) -- (5,9);
  \draw[bus] (4.6,9) -- (5.4,9);
  \draw (5,9) to[open, v^=$V_b^{(1)}$] (5,7.5);
  \draw[bus] (4.6,7.5) -- (5.4,7.5);
  \draw (5,7.5) -- (-0.8,7.5) -- (-0.8,6) -- (0,6);
  % seq. negativa
  \draw (0,6) to[L, l={$Z_2^{tot} = j0{,}1416$ pu}, i_=$I_b^{(2)}$] (4.3,6) -- (5,6);
  \draw[bus] (4.6,6) -- (5.4,6);
  \draw (5,6) to[open, v^=$V_b^{(2)}$] (5,4.5);
  \draw[bus] (4.6,4.5) -- (5.4,4.5);
  \draw (5,4.5) -- (-0.8,4.5) -- (-0.8,3) -- (0,3);
  % seq. zero (so a linha)
  \draw (0,3) to[L, l={$Z_{0,L} = j0{,}2666$ pu}, i_={$I_0 = 0$}] (4.3,3) -- (5,3);
  \draw[bus] (4.6,3) -- (5.4,3);
  \draw (5,3) to[open, v^=$V_b^{(0)}$] (5,1.5);
  \draw[bus] (4.6,1.5) -- (5.4,1.5);
  % retorno interrompido (Delta-Delta)
  \draw (5,1.5) -- (8,1.5) -- (8,1.2); \fill (8,1.2) circle (1.6pt);
  \draw (8,1.2) -- (8.4,0.8); \fill (8,0.5) circle (1.6pt);
  \draw (8,0.5) -- (8,0) -- (-3,0) -- (-3,7);
  \node[font=\scriptsize, align=left, anchor=west] at (8.6,1.0) {seq. 0 sem retorno\\($\Delta$-$\Delta$): $I = 0$};
  \node[draw, circle, inner sep=1pt] at (6.2,9.5) {F};
  \node[font=\small, align=center] at (2.5,-1.1) {Falta fase-terra: redes em série com $Z_0 = \infty$ $\Rightarrow$ $I_0 = I_1 = I_2 = 0$ $\Rightarrow$ $I_b^{(1\phi)} = 0$};
\end{circuitikz}
\end{document}
